\section{An effective frontier at a fixed persistent budget}\label{sec:effective}
The exact allocation invariant retains an entire law-valued trajectory. We now prove that, for an effectively presented class of positive instruments, its value admits certified finite approximations without increasing the observer's persistent alphabet. No stability of the normalized filter is assumed. The assertion is stronger than pointwise approximation of a fixed controller: its lower bound is uniform over all admissible Borel controllers of the declared size.

\subsection{The experimental approximation certificate}
Let $V$ be a finite-dimensional real vector space with a closed generating pointed cone $V_+$ and a strictly positive linear mass functional $e$. Assume the base $\mathcal B=\{x\in V_+:e(x)=1\}$ is compact. Use the base norm
\[
 \|v\|_{\mathcal B}=\inf_{v=v_+-v_-,\ v_\pm\in V_+}
                      \{e(v_+)+e(v_-)\}.
\]
A raw instrument $\mathsf I_t^a$ is a countably additive measure of positive linear maps on $V$, with $e\mathsf I_t^a(\mathcal Y_a)=e$. Preparation is $x_0\in\mathcal B$. The terminal loss is $L_u(x)$, where $0\le L_u\le He$ in the dual order. The compact decision set is $U$. Classical substochastic matrices and finite-dimensional completely positive quantum instruments are two instances, with $\ell^1$ and trace norm respectively. The physical state is not a register supplied to the observer.

Finite words of unnormalized instruments and terminal effects generate a finite-dimensional space of executable-test expectations. Restricting physical preparations to that space gives the compact affine predictive quotient. Weak continuity of the post-report law also follows in this finite-dimensional setting. Choose an interior normalized positive vector $x_*$; compactness of the base gives $x\preceq Cx_*$ for every normalized $x$. The report measure at $x_*$ dominates all vector instrument entries. For a continuous bounded function of the outgoing predictive state, its integrand is the report density times that function. It is continuous in $x$ wherever the density is positive, extends continuously by zero when the density vanishes, and is dominated by $C$ times the reference density. Dominated convergence proves continuity on the physical base. Representative independence and the compact quotient map descend it to the predictive carrier. Thus the exact theorem applies; the effective argument below uses the unnormalized instrument identities directly. A redundant physical representation may be used for offline calculation. Its dimension is not substituted for the dimension of an acquired law.

\begin{definition}[Support-resolved effective presentation]\label{def:effective}
For each refinement index $h$, each action has a finite Borel report partition $\mathcal C_a^h=(C_{a,b})_{b=1}^{B_a}$ and probability measures $\mathsf R_{a,b}$ supported on their cells. The partition and reconstruction measures are independent of phase and physical state. They come with a finite support table $S_{t,a}\subseteq\{1,\ldots,B_a\}$ such that, for every $x\in\mathcal B$,
\[
 e\mathsf I_t^a(\cdot)x\big|_{C_{a,b}}
 \quad\hbox{is equivalent to }\mathsf R_{a,b}\ \ (b\in S_{t,a}),
 \qquad \mathsf I_t^a(C_{a,b})=0\ \ (b\notin S_{t,a}).
\]
Equivalence concerns null sets, not a uniform likelihood-ratio bound. Cells inactive at every legal phase may be removed. Define the regenerated instrument
\[
 \widehat{\mathsf I}_t^a(dy)=
       \sum_b\mathsf R_{a,b}(dy)\mathsf I_t^a(C_{a,b}).
\]
There are supplied computable bounds $d_{t,h}\ge0$ such that
\begin{equation}\label{eq:instrument_error}
 \sup_{a\in\mathcal A_t,\,x\in\mathcal B}
 \int\big\|\big(\mathsf I_t^a-\widehat{\mathsf I}_t^a\big)(dy)x\big\|_{\mathcal B}
 \le d_{t,h},\qquad \sum_{t<n}d_{t,h}\longrightarrow0.
\end{equation}
The integral is the total variation norm of a finite-dimensional vector measure. Matrix entries of each cell instrument, preparation, and terminal loss at a given decision are computable, with certified absolute-error enclosures. For every $\rho>0$ a finite, explicitly encoded set $U_\rho$ is supplied such that every $u\in U$ has $v\in U_\rho$ with
\[
 \sup_{x\in\mathcal B}|L_u(x)-L_v(x)|\le\rho.
\]
Cell membership is an executable operation of declared description, scratch and time cost; alternatively an implementation provides a separate probability-of-misbinning certificate and preserves the report support signature. Nothing in computability of integrals alone supplies a bounded-cost analogue comparator.
\end{definition}
The support table permits phase flags and other clock-revealing reports. It need not be constant in $t$. Its role is to preserve \emph{exact} legality and stopping while eliminating the continuum report variable from a stationary programme.

Put $A=|\mathcal A|$, $B_\Sigma=\sum_a B_a$, and, for a dyadic denominator $Q=2^b$, set
\begin{equation}\label{eq:effective_constants}
 P=MA+M^2B_\Sigma,\qquad
 E_h=H\min\{1,\sum_{t<n}d_{t,h}\},\qquad
 \beta_Q=H\min\{1,n(A+M)/Q\}+\rho.
\end{equation}
Let $\mathscr G_{h,Q,\rho}$ contain all $M$-label autonomous programmes whose action and cell-to-label rows have denominator $Q$ and whose readouts lie in $U_\rho$, retaining only the programmes accepted by the finite support-graph check below. Write $K_\rho=|U_\rho|$. Unused labels and arbitrary rows at unreachable labels are allowed.

\begin{theorem}[Uniform finite certificates for the exact frontier]\label{thm:effective}
Fix $n\ge1$ and $M\ge1$ in a support-resolved effective experiment. Admissibility of each member of $\mathscr G_{h,Q,\rho}$ is decidable by finite graph reachability. The Borel autonomous class is empty if and only if this finite set is empty, for any supplied partition, any dyadic $Q\ge1$, and nonempty $U_\rho$. Moreover
\begin{equation}\label{eq:programme_count}
 |\mathscr G_{h,Q,\rho}|\le
 M2^M K_\rho^M(Q+1)^P.
\end{equation}
For each retained programme $g$, compute an interval $[l_g,u_g]$ containing its actual risk, of length at most $\eta>0$. If the class is nonempty, put $l=\min_g l_g$ and $u=\min_g u_g$. Then
\begin{equation}\label{eq:effective_interval}
 \boxed{\quad l-E_h-\beta_Q\le R_{\aut}(n,M)\le u,\qquad
 u-l\le\eta.\quad}
\end{equation}
A programme attaining $\min_g u_g$ is implementable on the original experiment at the same $M$, with risk at most $u$ and therefore excess over the optimal $M$-label Borel risk at most $E_h+\beta_Q+\eta$. The bounds may be clipped to $[0,H]$. They give an effective convergent two-sided approximation of $B_n^*+\Phi_{\aut}(n,M)$.

For an external width profile the same assertion holds, with phase-indexed rows and
\begin{equation}\label{eq:effective_external}
 P_{\ext}=\sum_{t<n}\left(m_t|\mathcal A_t|
                  +m_tm_{t+1}\sum_{a\in\mathcal A_t}B_a\right),\quad m_0=1,
\end{equation}
and count bound $K_\rho^{m_n}(Q+1)^{P_{\ext}}$. In $\beta_Q$ replace $n(A+M)$ by $\sum_{t<n}(|\mathcal A_t|+m_{t+1})$. An externally supplied phase remains free only in this second model.
\end{theorem}

\subsection{Proof: common reconstruction before discretizing programmes}
\begin{lemma}[Report reduction preserves one programme]\label{lem:report_reduction}
Let $R_h$ be the optimum when the reports are replaced by their cell indices, keeping the physical cell instruments $\mathsf I_t^a(C_{a,b})$ and the same observer interface. Under Definition~\ref{def:effective},
\begin{equation}\label{eq:report_reduction}
 R\le R_h\le R+E_h.
\end{equation}
Both transformations preserve the entire persistent alphabet, all positive-probability action legality constraints, and the exact deadline.
\end{lemma}
\begin{proof}
Any cell programme runs on the original reports by computing their cells; its joint physical-state/label law is exactly that given by the cell instruments. This proves $R\le R_h$.

Conversely, let $k(j\mid z,a,y)$ be any admissible stationary Borel programme. Define the single cell row
\[
 \bar k(j\mid z,a,b)=\int k(j\mid z,a,y)\,\mathsf R_{a,b}(dy).
\]
This is independent of phase. Running $k$ against the regenerated instrument has exactly the same state/label law as running $\bar k$ against the cell instrument. The reconstruction is an offline averaging identity, not a persistent report simulator or an exact posterior computation.

For the error, at any cut the joint state/label law is a family $(v_z)$ of positive vectors with $\sum_z e(v_z)=1$. A positive instrument followed by an action row and a report-to-label row is a contraction on the direct-sum base norm. Indeed decompose each signed input into positive parts; positivity and preservation of total mass bound the sum of output norms by the sum of input norms. On positive normalized inputs, replacing the instrument changes that sum by at most $d_{t,h}$, by \eqref{eq:instrument_error} and the stochastic row sums. Successive replacements therefore change the final joint state/label vector by at most $\sum_t d_{t,h}$. A terminal effect in $[0,He]$ changes its expected loss by at most $H$ times this norm, also at most $H$. This proves the risk error $E_h$. The argument is valid for quantum instruments without postulating a hidden classical trajectory for noncommuting states.

It remains to check the exact interface. For an active cell, an edge $z\to j$ after action $a$ is possible under the fine programme precisely when $\int k(j\mid z,a,y)d\mathsf R_{a,b}>0$. This equivalence holds at every active phase and every incoming positive conditional physical state. Finite paths have positive probability exactly when their action edges and active-cell edges have positive probability, by induction. Thus the original programme and its cell average have the same possible label/action paths at each phase. Inactive cells cannot occur in either. Premature halt, non-halt at $n$, and illegal actions are consequently absent from one programme if and only if they are absent from the other. Taking infima proves \eqref{eq:report_reduction}, also for classes whose optimum is not attained. The proof for external rows is identical with $t$ retained as a permitted argument.
\end{proof}

\begin{lemma}[Finite support graphs and dyadic approximation]\label{lem:dyadic_uniform}
For the cell experiment, a stationary programme is legal and stops exactly at $n$ if and only if the following finite recursion accepts it. Start with the initial label. At every $t<n$ reject a reachable stopping label or a positive action outside $\mathcal A_t$; otherwise follow all positive row entries for all $b\in S_{t,a}$. At $n$ require every reachable label to stop. For every admissible real-valued cell programme there is one in $\mathscr G_{h,Q,\rho}$ with risk increase at most $\beta_Q$.
\end{lemma}
\begin{proof}
The support equivalence and positivity induction in Lemma~\ref{lem:report_reduction} prove the graph criterion. No infinite-horizon recurrence is being tested. For a stochastic row of length $r$, take the floors of $Q$ times its entries and distribute the remaining units only to entries in its original support. The resulting row has the same or smaller support, total mass one, and total variation at most $r/Q$. This construction works even at $Q=1$; it does not require a lower bound on nonzero row probabilities. Apply it once to each shared stationary row, not separately at each phase. Removing edges cannot introduce a previously forbidden finite path, so legality and exact halting are preserved. There remain outgoing edges at every nonterminal row because its mass is one.

Couple the action and next-label choices up to their first discrepancy. At one call their total mismatch probability is at most $(A+M)/Q$. Conditional on equality so far, the physical experiment and label histories are the same. A union bound over $n$ calls bounds the loss increase by $H\min(1,n(A+M)/Q)$. Selecting a loss-$\rho$-net readout at each stopping label adds at most $\rho$. The same proof gives the displayed external sum. This comparison is uniform over the original row probabilities, including arbitrarily small positive entries.
\end{proof}

\begin{proof}[Proof of Theorem~\ref{thm:effective}]
Enumerate initial labels, stopping subsets, readouts, and numerator rows. There are at most $M$, $2^M$, and $K_\rho^M$ choices for the first three. A length-$r$ denominator-$Q$ row has at most $(Q+1)^r$ possibilities; summing the row lengths gives $P$. Apply the support check to each. If a Borel observer exists, Lemma~\ref{lem:report_reduction} and then Lemma~\ref{lem:dyadic_uniform} give a member of this finite list. The converse follows by execution, proving the empty-class assertion.

Let $R_G$ be the least actual risk in the retained finite list. The preceding lemmas give
\[
 R\le R_G\le R+E_h+\beta_Q.
\]
Each listed programme's risk is a finite sum of compositions of computable cell instruments, row probabilities and terminal effects. Certified interval arithmetic computes it to any prescribed absolute accuracy; computability here includes the oracle moduli in Definition~\ref{def:effective}, not just existence of the integrals. Hence $l\le R_G\le u$. If $i$ minimizes $l_i$, then $u\le u_i\le l_i+\eta=l+\eta$, proving the interval width. A programme minimizing $u_g$ has risk at most $u$. The error in \eqref{eq:effective_interval} follows. All steps for the external interface use the stated phase widths. Finally choose the computable refinement, $Q$, loss net and evaluation tolerance to make $E_h$, $\beta_Q$ and $\eta$ tend to zero. This is an effective value approximation, not an assertion that a general optimal Borel programme is attained or has a finite exact description.
\end{proof}

\begin{corollary}[Finite global obstruction and computation of excess]\label{cor:effective_obstruction}
For the effective class and any strict target $r<R_{\aut}(n,M)$, refinement eventually produces $l-E_h-\beta_Q>r$. The finite list of admissible programmes, its certified risk intervals, the support graph and the displayed error bounds form a global obstruction with candidate count bounded by \eqref{eq:programme_count}. This is not a bound on the number of affine tests in an arbitrary enumeration in Theorem~\ref{thm:finite_obstruction}.

The same construction on finite cell-history trees computes $B_n^*$: after report reduction there are finitely many action/report histories of depth at most $n$, and deterministic contingent actions and loss-net readouts suffice. Thus $\Phi_{\aut}$ and $\Phi_{\ext}$ have certified intervals obtained by subtracting a baseline interval for the \emph{same} experiment. The policy-dependent acquisition and compression terms in \eqref{eq:checkpoint} are not minimized separately.
\end{corollary}
\begin{proof}
Choose total interval error smaller than $R-r$. For the baseline, report reduction applies to full-history controllers with the external history retained; the coarsened history alphabet is finite. At a finite decision tree each action randomization can be purified backwards because the conditional objective is linear in that row. Replace terminal decisions by $U_\rho$. Finite enumeration and the same error calculation yield a baseline interval. Subtraction gives $[R_L-B_U,R_U-B_L]$, with the lower clipped at zero for the excess. This last computation may use vastly larger offline history trees; it does not grant extra state to the deployed $M$-label observer.
\end{proof}

\subsection{Programme, numerical and physical resources}
A listed programme has a direct read-only encoding of length
\begin{equation}\label{eq:effective_bits}
 O\big(P\log(Q+1)+M\log(K_\rho+1)+M+\log(M+1)\big)+C_{\rm cell}+C_{\rm out}.
\end{equation}
The supplied physical model, its integration algorithms, and a fixed interpreter are not included in this programme count. $C_{\rm cell}$ is the cell-classifier description and $C_{\rm out}$ encodes the finite decision net. They are not zero by convention. A dyadic row uses exactly $b$ fresh bits and a cumulative row scan; no rejection loop or retained random seed is required. In addition to the $\lceil\log_2M\rceil$ persistent bits, transition scratch is
\[
 O\big(b+\log(M+1)+\log(A+1)+\log(1+\max_aB_a)\big)+W_{\rm cell}+W_{\rm out}.
\]
Apart from physical acquisitions, the two scans use $O(n(A+M)(b+1))$ bit operations, plus cell and readout costs. Offline search inspects at most \eqref{eq:programme_count} candidates, each with a finite support check and a certified risk evaluation. No polynomial-time or optimal programme-length claim follows.

If numerical cell evaluation changes a report cell with total probability at most $\chi$ along the ideal deployed law, and preserves its support signature, first-discrepancy comparison adds at most $H\min(1,\chi)$ to the attained upper risk. Signature preservation retains the exact deadline. A readout error contributes its own loss modulus. A common-task simulator of defect $\varepsilon$ adds $H\varepsilon$ and multiplies persistent state by its certified simulator alphabet, as in Proposition~\ref{prop:simulation}. These are separate costs; they do not alter the ideal lower bound by assuming that every allowed numerical tolerance is attained.
